% Active section fragment; included through body/ch05_dual_domain_reconstruction.tex.
\section{真实与仿真触觉差异分析}
\label{sec:ch05_real_sim_difference}

真实与仿真触觉具有相同数组维度，并不意味着能够直接用于同一操作表示。本文采用的观测按16个Patch组织，共712个标量通道，仿真阵列覆盖五指及手掌。其传感器原生raw读数来自MuJoCo native touch，对感知区域内的法向接触力累加\cite{mujoco2026xml}；仿真processed观测则经过死区、限幅、非线性响应、时间滤波及Patch内邻域扩散处理（实现参数名为crosstalk）。Real processed观测来自实机自身采集处理流程，因此两域processed均应视为处理后的数值响应，不能直接称为真实压力或共同力学单位。对这种观测进行迁移，需要检查保留了哪些接触信息，而不只是核对输出形状。

P37旧开发数据的Real Train/Validation为5153/1408窗，Sim为1824/406窗；P43当前主线中，Real Train/Validation分别为5153/780窗，Natural Sim Train/Validation分别为4049/905窗，均采用$4\times712$窗口。P37统计不外推至P43，后者规模仅描述当前数据划分，不代表差异复现。后续统计按相应版本分别报告，并以来源分组考察样本独立性。

令$X_{tj}$为第$t$帧第$j$通道响应，精确非零Mask定义为$M_{tj}=\mathbf1[X_{tj}>0]$。在旧训练集上，Real与未转换Sim的激活比例分别约17.60\%和13.31\%，说明数值支持范围存在差异。Mask表示保留每帧712个激活位置，可以反映Patch内部支持及区域组合，却不是物理接触真值。滤波残留等微弱正值也会计入激活。因此，Mask首先反映数值支持范围；跨域热图按固定通道索引比较，物理配准的前提在下一节统一界定。

幅值分析仅取$X>0$的响应。旧训练集汇总正值中位数，Real约为0.00500250，未转换Sim约为$1.03974\times10^{-5}$；P95分别约1.50101与0.459485。这些是各域全部正值池的统计，不是逐Patch分位数的平均，也不是接触力误差。幅值分解以每帧正值的P10、P25、P50、P75、P90构成有效$4\times5$表示，其余通道仅补零以复用网络接口；不输入Taxel位置、正值数量或空Patch标志。对包含至少一帧全零响应帧的窗口，按照预定规则予以排除，故该表示只刻画有响应条件下的幅值，不能覆盖完全无接触状态。

表\ref{tab:ch05_domain_feature_difference}列出旧P37数据上既有转换输出经幅值校准后的表示分解诊断。Full与分位数先作log1p，各表示分别用Train拟合标准化，Probe由Logistic、kNN各1个及MLP、GRU、Transformer各Seed11、12、13组成，共11个；Full复用P38在相同输出、划分和预处理上的Probe拟合，不新增训练。此处为迁移后Probe诊断，非后文迁移前冻结的域分类评价基线（C0）验收。定义$D=1-2|\mathrm{BA}-0.5|$，只度量BA偏离0.5的程度；表中三个表示的最小D均低于0.90，表明至少存在Probe未满足固定阈值下的BA等价门槛；其排序判别能力仍需结合AUC判断。因此，不能仅凭这些结果认定支持或正值摘要已经实现域混淆。幅值调整保持Mask不变，不代表完整触觉已匹配。

\begin{table}[htbp]
 \centering\small
 \caption{旧P37数据校准后触觉的表示分解诊断}
 \label{tab:ch05_domain_feature_difference}
 \begin{tabular}{llc}
 \hline
 表示&有效内容&11个Probe实例中的最小D\\\hline
 Mask&$4\times712$精确非零支持&0.0272\\
 Amplitude&$4\times5$正值分位数&0.1494\\
 Full&$4\times712$完整响应&0.0573\\\hline
 \end{tabular}
\end{table}

三种表示的特征空间、预处理及纳入条件不同，D值不能相减计算信息贡献，也不证明支持与幅度独立。低维分位数舍弃了响应所在通道及跨Patch组合，即使边际分布相似，也可能无法维持原接触结构。差异还可能来自任务构成、采集协议与处理链，不能全部归因于传感器物理性能。

四帧窗口还承载激活持续、出现、消失及幅值变化。P37验证集第1至第2帧的Mask变化比例，Real约1.7154\%，未转换Sim约0.0958\%，分母为各自窗口与通道总数，而非条件转移概率。真实帧间隔并非恒定，帧间差分不能直接解释为物理时间导数或频率差异。幅值还需检查相邻帧全通道绝对差，不能仅用Mask持续性代替响应强度变化。因此，完整重建应同时保留激活支持、正值幅度及短时序关系。下一节据此在缺少同步物理配对的条件下定义跨域重建问题。
